%\chapter{作者简历及攻读学位期间发表的学术论文与研究成果}
%
%\textbf{本科生无需此部分}。
%
%\section*{作者简历}
%
%\subsection*{casthesis作者}
%
%吴凌云，福建省屏南县人，中国科学院数学与系统科学研究院博士研究生。
%
%\subsection*{ucasthesis作者}
%
%莫晃锐，湖南省湘潭县人，中国科学院力学研究所硕士研究生。
%
%\section*{已发表(或正式接受)的学术论文:}
%
%[1] ucasthesis: A LaTeX Thesis Template for the University of Chinese Academy of Sciences, 2014.
%
%\section*{申请或已获得的专利:}
%
%(无专利时此项不必列出)
%
%\section*{参加的研究项目及获奖情况:}
%
%可以随意添加新的条目或是结构。

\chapter[致谢]{致\quad 谢}\chaptermark{致\quad 谢}% syntax: \chapter[目录]{标题}\chaptermark{页眉}
%\thispagestyle{noheaderstyle}% 如果需要移除当前页的页眉
%\pagestyle{noheaderstyle}% 如果需要移除整章的页眉
\pagestyle{mainmatterstyle} % 与前文页眉页脚格式相同

感谢指导教师在论文选题、方法设计和写作过程中给予的悉心指导与支持。感谢上海科技大学信息科学与技术学院提供的计算资源（idata2 服务器），使本文的 PINN 模型训练得以顺利完成。感谢 ucasthesis 模板的开发者和 gbt7714-bibtex-style 的维护者，为本文的排版提供了便利。本文的 PINN 实现参考了 Dhiman \& Hu（2023）和 Guo 等人（2026）的开源工作，在此一并致谢。

\cleardoublepage[plain]% 让文档总是结束于偶数页，可根据需要设定页眉页脚样式，如 [noheaderstyle]

